% !TEX TS-program = pdflatex
% !TeX program = pdflatex
% !TEX encoding = UTF-8
% !TEX spellcheck = fr

\documentclass[11pt, a4paper]{article}
%\usepackage{fullpage}
\usepackage[left=1cm,right=1cm,top=1cm,bottom=2cm]{geometry}
\usepackage[fleqn]{amsmath}
\usepackage{amssymb}
%\usepackage{indentfirst}
\usepackage[T1]{fontenc}
\usepackage[utf8]{inputenc}
\usepackage[french,english]{babel}
\usepackage{txfonts} 
\usepackage[]{graphicx}
\usepackage{multirow}
\usepackage{hyperref}
\usepackage{parskip}
\usepackage{multicol}
\usepackage{wrapfig}
\usepackage{fancyhdr}
\usepackage{hyperref}
\usepackage{tcolorbox}
\usepackage{natbib}

\usepackage{turnstile}%Induction symbole

\renewcommand{\baselinestretch}{1}

\setlength{\parindent}{24pt}

%\usepackage{etoolbox}
%\patchcmd{\thebibliography}{\section*{\refname}}{}{}{}

\fancyhf{}

\lfoot{ARIES Abdelkrime}
\cfoot{ESI/TALN}
\rfoot{\textbf{\thepage}}

\renewcommand{\headrulewidth}{0pt} % remove lines as well
\renewcommand{\footrulewidth}{1pt}

\newcommand\repeatstr[1]{\leavevmode\xleaders\hbox{#1}\hfill\kern0pt}

\renewcommand{\bibsection}{}
\bibliographystyle{humannat}

\begin{document}

\selectlanguage {french}
%\pagestyle{empty} 
\pagestyle{fancy}

\noindent
\begin{tabular}{ll}
\multirow{3}{*}{\includegraphics[width=2cm]{esi-logo.png}} & \'Ecole national Supérieure d'Informatique\\
& 2ème année cycle supérieure\\
& Traitement automatique du langage naturel (TALN)
\end{tabular}\\[.25cm]
\noindent\rule{\textwidth}{1pt}\\%[-0.25cm]
\begin{center}
{\LARGE \textbf{Programme}}
\begin{flushright}
	ARIES Abdelkrime
\end{flushright}
\end{center}
\noindent\rule{\textwidth}{1pt}

\begin{abstract}
	Le traitement automatique du langage naturel (TALN), aussi connu comme "traitement automatique des langues" (TAL), est un domaine qui implique la linguistique et l'informatique (et l'intelligence artificiel pour ceux qui séparent entre l'IA et l'informatique). Il fait partie des domaines de l'intelligence artificiel.
	
	Récemment, TALN commence à recevoir plus d'attention à cause de ces applications surtout dans les entreprises. Parmi les bénéfices du TALN dans le domaine d'affaires, on peut citer : 
	\begin{itemize}
		\item Augmenter la productivité en utilisant des applications comme la traduction automatique et le résumé automatique (pourtant ces deux applications sont loin d'être parfaites)
		\item Service Clientèle : la réponse automatique aux questions des clients en utilisant les chatbots (question-réponse et reconnaissance de voix). 
		\item Surveillance de la réputation : on utilise l'analyse des sentiments pour savoir si les clients sont heureux avec ses produits ou non. 
		\item La publicité : on scannant les réseaux sociaux et les courriels, on peut savoir qui est intéressé par ses produits. Ceci permet aux entreprises de viser l'audience de la publicité. 
		\item Connaissance du marché (Market intelligence) : surveiller les compétiteurs afin de se tenir au courant des évènements liés à l'industrie.
		
	\end{itemize}
\end{abstract}

\section{Informations générales}

\begin{minipage}{0.49\textwidth}

\subsection{Pré-requis}

\begin{itemize}
	\item Théorie des langages de programmation et applications (THP)
	\item Probabilités et statistiques
	\item Algèbre et analyse 
	\item Logique des prédicats du premier ordre
	\item Apprentissage automatique 
	\item Programmation (surtout en Python et Java)
\end{itemize}
\end{minipage}
\begin{minipage}{0.49\textwidth}


\subsection{Objectifs}

\begin{itemize}
	\item Appliquer des notions mathématiques apprises au long du cursus sur des problèmes réels de la langue 
	\item Apprendre quelques notions linguistiques et un peu de philosophie (représentation de connaissance) 
	\item Découvrir quelques outils et ressources du TALN, et programmer des solutions simples à quelques problèmes
	\item Savoir appliquer les concepts vus en THP sur des différentes langues (langages naturels)
%	\item Avoir une idée sur quelques travaux TALN réalisés au sein de l'ESI
\end{itemize}
\end{minipage}

\section{Plan du module}

Les estimations horaires sont optimistes ; les cours peuvent prendre plus du temps que prévu. 
Les derniers chapitres (concernant les applications du TALN) sont optionnels ; ils seront présentés selon la disponibilité temps. On estime 30 heurs du cours, ce qui fait 15 séances de 2h. 

\begin{itemize}
	\item Chapitre 1. Introduction (1h)
	\item Chapitre 2. Traitements basiques du texte (4h)
	\item Chapitre 3. Les modèles de langages (2h)
	\item Chapitre 4. Étiquetage morpho-syntaxique (2h)
	\item Chapitre 5. Analyse syntaxique (4h)
	\item Chapitre 6. Sens des mots et désambiguïsation lexicale (4h)
	\item Chapitre 7. Analyse sémantique (4h)
	\item Chapitre 8. Détection de la coréférence (2h)
	\item Chapitre 9. La cohérence du discours (2h)
	\item Chapitre 10. Quelques applications (5h)
\end{itemize}

%\clearpage

\section{Plan détaillé des cours}

\begin{tcolorbox}

\subsection*{Chapitre 1. Introduction} %\repeatstr{.}

\begin{multicols}{2}
\begin{itemize}
	\item Histoire
	\begin{itemize}
		\item Naissance de l'IA et âge d'or
		\item Hiver de l'IA
		\item Printemps de l'IA
	\end{itemize}
	\item Les niveaux de traitement d'une langue
	\begin{itemize}
		\item Phonétique, phonologie et orthographe
		\item Morphologie et syntaxe
		\item Sémantique
		\item Pragmatique et discours
	\end{itemize}
	\item Applications du TALN
	\begin{itemize}
		\item Tâches
		\item Systèmes
		\item Affaires (business)
	\end{itemize}
	\item Défis du TALN
	\begin{itemize}
		\item Ressources
		\item Compréhension de la langue
		\item \'Evaluation
		\item \'Ethique
	\end{itemize}
\end{itemize}
\end{multicols}
\end{tcolorbox}

\subsection*{\uppercase{Morphologie et lexique}}

\begin{tcolorbox}
\subsubsection*{Chapitre 2. Traitements basiques du texte}

\begin{multicols}{2}
\begin{itemize}
	\item Caractères
	\begin{itemize}
		\item Expressions régulières
		\item Distance d'édition
	\end{itemize}
	\item Segmentation du texte
	\begin{itemize}
		\item Délimitation de la phrase
		\item Séparation des mots
	\end{itemize}
	\item Normalisation et filtrage du texte
	\begin{itemize}
		\item Normalisation du texte
		\item Filtrage du texte
	\end{itemize}
	\item Morphologie
	\begin{itemize}
		\item Formation des mots
		\item Réduction des formes
	\end{itemize}
\end{itemize}
\end{multicols}
\end{tcolorbox}

\subsection*{\uppercase{Syntaxe}}

\begin{tcolorbox}
\subsubsection*{Chapitre 3. Modèles de langues}
	
\begin{multicols}{2}
	\begin{itemize}	
		\item Modèle N-gramme
		\begin{itemize}
			\item Formulation
			\item Lissage (Smoothing)
		\end{itemize}
		\item Modèles neuronaux
		\begin{itemize}
			\item Réseau de neurones à propagation avant
			\item Réseau de neurones récurrents
			\item Quelques améliorations
		\end{itemize}
		\item Évaluation
		\begin{itemize}
			\item Approches
			\item Perplexité
		\end{itemize}
	\end{itemize}
\end{multicols}
\end{tcolorbox}

\begin{tcolorbox}
\subsubsection*{Chapitre 4. Étiquetage morpho-syntaxique}

\begin{multicols}{2}
\begin{itemize}
	\item \'Etiquetage de séquences
	\begin{itemize}
		\item Description
		\item Applications
	\end{itemize}
	\item Description de la tâche
	\begin{itemize}
		\item Classes universelles
		\item Treebanks
		\item Difficulté et outils
	\end{itemize}
	\item Approches
	\begin{itemize}
		\item Modèle de Markov
		\item Maximum entropy
		\item Modèle neuronal
	\end{itemize}
\end{itemize}
\end{multicols}
\end{tcolorbox}

\begin{tcolorbox}
\subsubsection*{Chapitre 5. Analyse syntaxique}

\begin{multicols}{2}
	\begin{itemize}
		
		\item Structures syntaxiques
		\begin{itemize}
			\item Annotation constituante
			\item Annotation fonctionnelle
		\end{itemize}
		
		\item Analyse des constituants
		\begin{itemize}
			\item Algorithme CKY
			\item Algorithme CKY probabiliste
		\end{itemize}
	
		\item Analyse des dépendances
		\begin{itemize}
			\item Par transition
			\item Par graphe
		\end{itemize}
		
	\end{itemize}
\end{multicols}
\end{tcolorbox}

\subsection*{\uppercase{Sémantique}}

\begin{tcolorbox}
\subsubsection*{Chapitre 6. Sématique lexicale}

\begin{multicols}{2}
\begin{itemize}
	
	\item Bases de données lexicales
	\begin{itemize}
		\item Relations sémantiques
		\item Wordnet
		\item Autres ressources
	\end{itemize}
	
	\item Représentation vectorielle des mots
	\begin{itemize}
		\item TF-IDF
		\item Mot-Mot
		\item Analyse sémantique latente (LSA)
	\end{itemize}

	\item Word embedding
	\begin{itemize}
		\item word2vec
		\item GloVe
		\item Embeddings contextuels
		\item Évaluation des modèles
	\end{itemize}

	\item Désambiguïsation lexicale
	\begin{itemize}
		\item Basée sur des bases de connaissance
		\item Basée sur l'apprentissage automatique
	\end{itemize}
	
\end{itemize}
\end{multicols}
\end{tcolorbox}

\begin{tcolorbox}
\subsubsection*{Chapitre 7. Sémantique de phrases}

\begin{multicols}{2}
	\begin{itemize}
		
		\item Rôles sémantiques
		\begin{itemize}
			\item Rôles thématiques
			\item FrameNet 
			\item PropBank	
		\end{itemize}
	
		%https://medium.com/explorations-in-language-and-learning/trends-in-semantic-parsing-part-1-ba11888523cb
		\item Étiquetage de rôles sémantiques (Semantic role labeling)
		\begin{itemize}
			\item En utilisant des caractéristiques
			\item En utilisant les réseaux de neurones
		\end{itemize}
		
		\item Représentation sémantique des phrases
		\begin{itemize}
			\item Logique du premier ordre
			\item Graphes (AMR)
		\end{itemize}
		
	\end{itemize}
\end{multicols}
\end{tcolorbox}

\subsection*{\uppercase{Énonciation et pragmatique}}

\begin{tcolorbox}
\subsubsection*{Chapitre 8. Détection de la coréférence}

\begin{multicols}{2}
	\begin{itemize}
		
		\item Références
		\begin{itemize}
			\item Formes des références 
			\item Manière de référencement
			\item Propriétés des relations de coréférence
		\end{itemize}
	
		\item Résolution des coréférences
		\begin{itemize}
			\item Détection de mention
			\item Liaison
			\item Évaluation
		\end{itemize}
	
		\item Tâches connexes
		\begin{itemize}
			\item Annotation sémantique (Entity linking)
			\item Reconnaissance des entités nommées
		\end{itemize}
		
	\end{itemize}
\end{multicols}
\end{tcolorbox}

\begin{tcolorbox}
\subsubsection*{Chapitre 9. Cohérence du discours}

\begin{multicols}{2}
	\begin{itemize}
	
		\item Relations de cohérence
		\begin{itemize}
			\item Rhetorical Structure Theory (RST)
			\item Penn Discourse TreeBank (PDTB)
		\end{itemize}
	
		\item Analyse basée structure de discours
		\begin{enumerate}
			\item Analyse RST
			\item Analyse PDTB
		\end{enumerate}
	
		\item Analyse basée sur l'entité de discours
		\begin{enumerate}
			\item Centering theory
			\item Entity Grid model
		\end{enumerate}
		
	\end{itemize}
\end{multicols}
\end{tcolorbox}

\subsection*{\uppercase{Quelques applications}}

\begin{tcolorbox}
\subsubsection*{Chapitre 10. Quelques applications}

\begin{multicols}{2}
	\begin{itemize}
		\item Transformation
		\begin{itemize}
			\item Traduction automatique de textes
			\item Résumé automatique de textes
		\end{itemize}
		\item Interaction
		\begin{itemize}
			\item Questions-Réponses
			\item Systèmes de dialogue et chatbots
		\end{itemize}
	    \item Classification
	    \begin{itemize}
	    	\item Analyse des sentiments
	    	\item Lisibilité
	    \end{itemize}
    	\item Voix
		\begin{itemize}
			\item Reconnaissance vocale
			\item Synthèse vocale
		\end{itemize}
	\end{itemize}
\end{multicols}
\end{tcolorbox}


\section{Plan des TPs}

On choisit seulement une ou deux applications par chapitre, selon la difficulté et le nombre de séances.

\begin{itemize}
	
	\item Chapitre 2. Traitements basiques du texte
	\begin{itemize}
		\item Fouille des contacts en utilisant les expressions régulières
		\item Correction des fautes d'orthographe OU comparaison des séquences d'ADN
		\item Réalisation d'un programme de racinisation en utilisant Snowball
		\item Réalisation d'un petit moteur de recherche pour une langue aux choix (sans utiliser des APIs comme elasticsearch)
	\end{itemize}

	\item Chapitre 3. Modèles de langues
	\begin{itemize}
		\item Système d'auto-complétion
		\item Classification des documents (spam/non spam)
		\item Identification de la langue : fréquences des mot vs. NGrams (caractères) %http://www.practicalcryptography.com/miscellaneous/machine-learning/tutorial-automatic-language-identification-ngram-b/
	\end{itemize}

	\item Chapitre 4. Étiquetage morpho-syntaxique
	\begin{itemize}
		%https://github.com/amjha/HMM-POS-Tagger
		\item Réalisation d'un programme d'étiquetage morpho-syntaxique en utilisant les modèles de Markov cachés
	\end{itemize}

	\item Chapitre 5. Analyse syntaxique
	\begin{itemize}
		\item Implémentation de l'algorithme CKY et application sur une forme simplifiée d'une langue 
	\end{itemize}

	\item Chapitre 6. Sens des mots et désambiguïsation lexicale
	\begin{itemize}
		%https://addi.ehu.es/bitstream/handle/10810/23867/TESIS_GONZALEZ_AGIRRE_AITOR.pdf?sequence=1
		\item Réalisation d'un programme simple de détection de plagiat en utilisant la similarité entre les phrases. Trois variantes : lexicale (Similarité cosinus), sémantique (Wordnet), statistique (réseaux de neurones : encoder-decoder)
		\item Désambiguïsation lexicale (Wordnet, Réseaux de neurones)
	\end{itemize}

	\item Chapitre 7. Analyse sémantique
	\begin{itemize}
		\item Réalisation d'un système de question-réponse simple (réponse automatique)
	\end{itemize}

	\item Chapitre 8. Détection de la coréférence
	\begin{itemize}
		\item Amélioration du système de question-réponse en utilisant la détection de la coréférence
	\end{itemize}

	\item Chapitre 9. La cohérence du discours
	\begin{itemize}
		\item Ordonnancement des phrases
		\item Évaluation de la cohérence
	\end{itemize}

	Chapitre 10. Quelques applications
	\begin{itemize}
		\item Traduction automatique en utilisant un réseaux de neurones Seq2Seq avec Attention
		\item Résumé automatique (comparaison entre plusieurs méthodes)
		\item Analyse des sentiments 
		\item Chabot 
		\item Réalisation d'un système de reconnaissance de paroles (deep learning)
		\item Utilisation d'un API pour synthétiser le son (par exemple, pyttsx)
		\item Fusionner avec le système question réponse
	\end{itemize}

\end{itemize}

%\section{Outils}
%
%L'utilisation des outils dépend des TPs proposés. Cette liste contient 
%
%\begin{itemize}
%	\item  
%\end{itemize}

\section{Bibliographie}

\nocite{*}

%\bibliographystyle{apalike}

\bibliography{biblio}

\end{document}
